\documentclass[12pt]{article}
\usepackage{amsmath,amssymb,amsfonts}
\usepackage[margin=1in]{geometry}

\begin{document}

\textbf{Chapter 6, Problem 19.} \textit{Cauchy's integral formula for a polynomial.} Let $P(z)$ be a polynomial. Using Cauchy's theorem, prove that
\[
\frac{1}{2\pi i}\oint_C P(z)\,dz = P(a),
\]
where $C$ encloses the point $a$. [\textit{Hint:} Consider $(z-a)^n\frac{1}{z-a}$, which is also a polynomial (why?), and evaluate the integral you obtain from this expression.]

\bigskip
\textbf{Solution.}

We want to show:
\[
\frac{1}{2\pi i}\oint_C \frac{P(z)}{z-a}\,dz = P(a),
\]
where $C$ is a closed contour enclosing $a$, and $P(z)$ is a polynomial.

This is simply Cauchy's integral formula applied to $P(z)$, which is entire (analytic everywhere). Since $P(z)$ is analytic inside and on $C$, Cauchy's integral formula gives directly:
\[
P(a) = \frac{1}{2\pi i}\oint_C \frac{P(z)}{z-a}\,dz. \quad \checkmark
\]

\textbf{Alternative proof using the hint:} Write $P(z) = \sum_{k=0}^n c_k z^k$. Perform polynomial division of $P(z)$ by $(z-a)$:
\[
P(z) = (z-a)Q(z) + P(a),
\]
where $Q(z)$ is a polynomial of degree $n-1$ (this is the remainder theorem from algebra).

Therefore:
\[
\frac{P(z)}{z-a} = Q(z) + \frac{P(a)}{z-a}.
\]

Integrating:
\[
\oint_C \frac{P(z)}{z-a}\,dz = \oint_C Q(z)\,dz + P(a)\oint_C \frac{dz}{z-a}.
\]

Since $Q(z)$ is a polynomial (hence entire), $\oint_C Q(z)\,dz = 0$ by the Cauchy--Goursat theorem.

And $\oint_C \frac{dz}{z-a} = 2\pi i$ (the fundamental result of complex analysis).

Therefore:
\[
\oint_C \frac{P(z)}{z-a}\,dz = 0 + P(a)\cdot 2\pi i = 2\pi i\,P(a).
\]
\[
\boxed{\frac{1}{2\pi i}\oint_C \frac{P(z)}{z-a}\,dz = P(a).}
\]

\end{document}
